\documentclass[../main.tex]{subfiles}
\begin{document}
\chapter{Convergence of Random Variables}
\section{Convergence in Probability}
\begin{definition}[Convergence in Probability]
  A sequence of random variables $(X_n: n \in \N)$ \textit{converges to a random variable $X$ in probability} if:
  \[
    \forall \varepsilon > 0,\ \P(|X_n - X| > \varepsilon) \to 0 \text{ as } n \to \infty
  \]
  we then write:
  \[
    X_n \xrightarrow{P} X \text{ or } \underset{n\to\infty}{\operatorname{plim}} X_n = X
  \]
\end{definition}
\begin{theorem}[Weak Law of Large Numbers]
  Let $(X_n)$ be independent and identically distributed random variables with finite $\mu = \E[X_1]$.

  If we set $S_n = X_1 + \cdots + X_n$ for $n \in \N$, then:
  \[
    \frac{S_n}{n} \xrightarrow{P} \mu \text{ as } n \to \infty
  \]
  That is:
  \[
    \forall \varepsilon > 0,\ \P\left(\abs{\frac{S_n}{n} - \mu} > \varepsilon\right) \to 0 \text{ as } n \to \infty
  \]
\end{theorem}
\begin{remark}
  In the definition of convergence in probability, we specified that a sequence converges to \textbf{random variable} $X$, here $\mu$ is just a constant random variable so we can still apply the definition.
\end{remark}
\begin{proof}
  For ease, in this proof, we will assume $\sigma^2 = \Var(X_1) < \infty$.
  \begin{align*}
    \P\left(\abs{\frac{S_n}{n} - \mu} > \varepsilon\right) &= \P(|S_n - n\mu| > n\varepsilon) \\
                                                           &\leq \frac{\E[|S_n - n\mu|^2]}{n^2 \varepsilon^2} \text{ using Markov's -- \cref{markovs}} \\
                                                           &= \frac{\Var(S_n)}{n^2 \varepsilon^2} \text{ as $\E[S_n] = n\mu$}
  \end{align*}
  Since the $X_i$ are i.i.d, we have that $\Var(S_n) = n\sigma^2$ so:
  \[
    \P\left(\abs{\frac{S_n}{n} - \mu} > \varepsilon\right) \leq \frac{n \sigma^2}{n^2 \varepsilon^2} = \frac{\sigma^2}{n \varepsilon^2} \to 0 \text{ as } n \to \infty
  \]
  as required.
\end{proof}
\section{Convergence Almost Surely}
\begin{definition}
  Let $(X_n)_{n \in \N}$ be a sequence of random variables at let $X$ be another random variable.
  We say that \textit{$X_n$ converges to $X$ almost surely} or \textit{with probability 1} if:
  \[
    \P\left(\lim_{n \to 0} X_n = X\right) = 1
  \]
  We then write $X_n \to X$ as $n \to \infty$ \textit{almost surely}.
\end{definition}
\begin{remark}
  When we say:
  \[
    \P\left(\lim_{n \to \infty} X_n = X\right) = 1
  \]
  We mean that:
  \[
    \exists A \in \salg \text{ s.t. } \P(A) = 1 \text{ and } \forall \omega \in A\ \forall \varepsilon > 0 \ \exists n_0(\varepsilon, \omega) \text{ s.t. } \forall n\geq n_0\ |X_n(\omega) - X(\omega)| \leq \varepsilon
  \]
  So there is a set $A$ with probability one where the numerical sequence $(X_n(\omega))_n$ converges to $X(\omega)$ for all $\omega \in A$.
\end{remark}
\begin{proposition}[Convergence almost surely $\implies$ Convergence in probability]
  If $X_n \to 0$ as $n \to \infty$ almost surely, then $X_n \xrightarrow{P} 0$ as $n \to \infty$.
\end{proposition}
\begin{proof}
  We need to show that $\forall \varepsilon > 0\ \P(|X_n| > \varepsilon) \to 0$ as $n \to \infty$, or equivalently, $\forall \varepsilon > 0\ \P(|X_n| \leq \varepsilon) \to 1$ as $n \to \infty$.

  To show that $\P(|X_n| \leq \varepsilon) \to 1$, we want to bound this below by something that increases to 1 in the limit.
  We can start by bounding it as:
  \[
    \P(|X_n| \leq \varepsilon) \geq \P\left(\bigcap_{m = n}^{\infty} \{|X_m| \leq \varepsilon\}\right)
  \]
  since $\bigcap_{m = n}^{\infty} \{|X_m| \leq \varepsilon \} \subseteq \{|X_n| \leq \varepsilon \}$.
  Then let $A_n = \bigcap_{m = n}^{\infty} A_n$.
  This is an increasing sequence of events as $A_{n} \subseteq A_{n + 1}\ \forall n$ so by continuity of probability measures (\cref{continuityProbMeasures}):
  \begin{align*}
    \P(A_n) \to \P\left(\bigcup_{n = 0}^{\infty} A_n\right) &= \P\left(\bigcup_{n = 0}^{\infty} \bigcap_{m = n}^{\infty} \{|X_n| \leq \varepsilon\}\right) \\
                                                            &= \P(|X_m| \leq \varepsilon \text{ for all $m$ sufficiently large}) \\
                                                            &\geq \P(\forall \varepsilon > 0\ |X_m| \leq \varepsilon \text{ for all $m$ sufficiently large}) \\
                                                            &= 1 \text{ as $X_n \to 0$ almost surely}
  \end{align*}
  Thus $\P(|X_n| \leq \varepsilon) \to 1$ as $n \to \infty$ so $X_n \xrightarrow{P} 0$ as $n \to \infty$.
\end{proof}
\begin{theorem}[Strong Law of Large Numbers -- SLLN]
  Let $(X_n)$ be independent and identically distributed random variables with finite $\mu = \E[X_1]$.

  If we set $S_n = X_1 + \cdots + X_n$ for $n \in \N$, then:
  \[
    \frac{S_n}{n} \to \mu \text{ as } n \to \infty \text{ almost surely}
  \]
\end{theorem}
\begin{proof}
  \nonexaminable
  For ease, in this proof, we will assume that $\E[X^4_1] < \infty$.

  Set $Y_i = X_i - \mu$ so that $\E[Y_i] = 0$.
  We can then check that $\E[Y^4_i]$ is also finite:
  \[
    \E[Y^4_i] = \E[(X_1 - \mu)^4] \leq 2^{4} \cdot (\E[X^4_1] + \mu^4) < \infty
  \]
  Define $S_n = \sum_{i = 1}^{n} Y_i$, we need to show that $\frac{S_n}{n} \to 0$  as $n \to \infty$ almost surely.
  To do this, our goal is to show that the series $\sum_{n = 1}^{\infty} (\frac{S_n}{n})^4$ is finite with probability 1.
  By the $n$-th term test (See Analysis I), if a series converges, its terms must go to 0 and therefore if $\sum_{n = 1}^{\infty} (\frac{S_n}{n})^4$ is finite with probability 1, $\frac{S_n}{n} \to 0$ as $n \to \infty$ with probability 1.

  To show the sum is finite with probability 1, it suffices to show that its expectation is finite since this implies that it is finite with probability one.

  Since $(\frac{S_n}{n})^4$ are non-negative, by \cref{sumExpectationSwap}, we can use the linearity of expectation on the infinite sum:
  \begin{align*}
    \E\left[\sum_{n = 1}^{\infty} \left(\frac{S_n}{n}\right)^4\right] &= \sum_{n = 1}^{\infty}\frac{1}{n^4} \E[S^4_n]
  \end{align*}
  We therefore want to bound $\E[S^4_n]$:
  \begin{align*}
    \E[S^4_n] &= \E[(Y_1 + \cdots Y_n)^4] \\
              &= \E\left[\sum_{i = 1}^{n} Y^4_i + \binom{4}{2}\sum_{1 \leq i < j \leq n}^{n} Y^2_i Y^2_j + R\right]
  \end{align*}
  where $R$ is a sum of terms of the form $Y^3_iY_j$, $Y^2_iY_jY_k$ and $Y_iY_jY_kY_l$ with distinct indices $i, j, k, l$.
  Since the $Y_i$ are independent and $\E[Y_i] = 0$, we then have $\E[R] = 0$.
  Note also that the $\binom{4}{2}$ originates from the fact that, for each $i < j$, we pick 2 of the 4 brackets to pick the $Y_i$ terms from, and then pick the $Y_j$ terms from the other two.

  For the terms with two squares:
  \begin{align*}
    \E[Y^2_iY^2_j] &= (\E[Y^2_1])^2 \text{ as they are i.i.d}\\
                   &\leq \E[Y^4_1] < \infty \text{ using Jensen's inequality -- \cref{jensensInequality}}
  \end{align*}
  Therefore we can bound $\E[S^4_n]$ as:
  \begin{align*}
    \E[S^4_n] &\leq n \E[Y^4_1] + 6 \cdot \binom{n}{2} \cdot \E[Y^4_1] \text{ as there are $\binom{n}{2}$ ways to pick $i < j$} \\
              &\leq 3n^2 \E[Y^4_1]
  \end{align*}
  So:
  \[
    \E\left[\sum_{n = 1}^{\infty} \left(\frac{S_n}{n}\right)^4\right] = \sum_{n = 1}^{\infty} \frac{\E[S^4_n]}{n^4} \leq \sum_{n = 1}^{\infty} \frac{3}{n^2} \cdot \E[Y^4_1] < \infty
  \]
  as $\sum_{n = 1}^{\infty} \frac{1}{n^2}$ converges ans $\E[Y^4_1] < \infty$.
\end{proof}
\section{Central Limit Theorem}
\subsection{Prelude}
Let $X_1, \ldots, X_n$ be independent and identically distributed random variables with $\mu = \E[X_1] < \infty$ and $\sigma^2 = \Var(X_1) < \infty$.
Again, define $S_n = X_1 + \cdots X_n$.
By the strong law of large numbers above, we expect $S_n \approx n\mu$ for large $n$.

But how does $S_n$ fluctuate around $n \mu$?
We can calculate the variance and expectation of these fluctuations:
\[
  \Var(S_n - n\mu) = \Var(S_n) = n \sigma^2 \text{ and } \E[S_n - n\mu] = 0
\]
The random variable $\frac{S_n - n\mu}{\sqrt{n \sigma^2}}$ then has expectation 0 and variance 1.
We can rewrite this random variable as:
\[
  \frac{S_n/n - \mu}{\sqrt{\sigma^2/n}} = \frac{S_n/n - \mu}{\sqrt{\Var(S_n/n - \mu)}}
\]
since $\Var(S_n/n - \mu) = \sigma^2/n$.
So if we take $\frac{S_n}{n} - \mu$, which we know converges to 0 almost surely, and then divide it by the square-root of its variance, we get a random variable with expectation 0 and variance 1.

Consider the case where the $X_i$ are independent and identically distributed  as $\normal(\mu, \sigma^2)$.
Since the sum of Gaussians is also a Gaussian (see \cref{sumOfNormals}), we have $S_n - n\mu \sim \normal(0, n\sigma^2)$ and therefore $\frac{S_n - n\mu}{\sigma \sqrt{n}} \sim \normal(0, 1)$.
So in this case, it follows a standard normal.
Although we do not yet know what the distribution of this is when working with a general distribution for $X_i$, this has the same expectation and variance as when we were just using any distribution for the $X_i$.

As we will see, the normal distribution is \textit{universal} as it appears as the limit of the distribution $\frac{S_n - n\mu}{\sigma \sqrt{n}}$, regardless of the distribution of $X_i$.
This is concept of the \textit{central limit theorem}.
\subsection{Convergence in Distribution}
\begin{definition}[Convergence in Distribution]
  A sequence of random variables $(X_n)_{n \in N}$ \textit{converges to a random variable $X$ in distribution} as $n \to \infty$ if:
  \[
    F_{X_n}(x) \to F_X(x) \text{ as $n \to \infty$}\ \forall x \in \R \text{ that are continuity points of $F_X$}
  \]
  where $F_{X_n}(x) = \P(X_n \leq x)$ is the probability distribution function.
  We then write $X_n \xrightarrow{(d)} X$ as $n \to \infty$.
\end{definition}
\begin{remark}
  In \cref{pdfProperties}, we saw that left limits always exist for $F$ and $F$ is always right continuous.
  This means that the continuity points are just the points where the left limit equals the right limit (which, by continuity, is the value of the function at that point).
\end{remark}
\end{document}
